\documentclass[11pt]{article}

\usepackage[margin=1.5cm]{geometry}
\usepackage{paracol}
\usepackage{hyperref}
\hypersetup{
    colorlinks=true,
    linkcolor=black,
    urlcolor=black,
    citecolor=black
}\usepackage{setspace}

\setlength{\parindent}{0pt}
\setlength{\parskip}{4pt}

\begin{document}

% ================= HEADER =================

\begin{paracol}{2}
\switchcolumn[0]


{\textbf{Florent TURRI}}

\vspace{2mm}

\switchcolumn

\begin{flushright}
Île-de-France, France \\
+33 6 69 53 73 32 \\
\textbf{\href{mailto:florentturri@gmail.com}{florentturri@gmail.com} \\}
\textbf{\href{https://www.linkedin.com/in/florent-turri}{www.linkedin.com/in/florent-turri} \\}
Permis A, B
\end{flushright}

\end{paracol}

\begin{center}
{\textbf{INGÉNIEUR LOGICIEL EMBARQUÉ CONFIRMÉ / SENIOR – C / C++ / LINUX}}

% ================= PROFILE =================

\textbf{PROFIL}

Ingénieur logiciel embarqué avec \textbf{6 années d’expérience en développement C/C++ sous Linux} dans des environnements industriels critiques, notamment pour des \textbf{grands comptes du secteur défense et systèmes complexes}.

Expert en développement bas niveau, intégration logiciel/système, debugging et optimisation sur plateformes Linux embarquées. Forte maîtrise des environnements Unix depuis plus de 10 ans.
\end{center}
\vspace{5mm}

% ================= COLUMNS =================

\columnratio{0.4,0.6}
\begin{paracol}{2}
\switchcolumn[0]
% ================= LEFT COLUMN =================

\underline{\textbf{COMPÉTENCES CLÉS}}\\
\textbf{Langages}

C \\
C++ \\
shell \\
elisp \\
Java \\
python \\
sql

\vspace{3mm}

\textbf{Frameworks / librairies}

Qt \\
GLib \\
boost \\
MinGW

\vspace{3mm}

\underline{\textbf{ENVIRONNEMENTS}}\\
\underline{\textbf{DE DÉVELOPPEMENT}}\\
Linux \\
Windows

\vspace{3mm}

\textbf{SYSTÈMES CIBLES}

\textbf{Logiciels} \\
Linux embarqué \\
IBMi / AS400

\vspace{2mm}

\textbf{Matériels} \\
Centrale inertielle \\
Systèmes embarqués (Arduino, STM32) \\
FPGA

\vspace{3mm}

\underline{\textbf{CONCEPTS CLÉS}}\\
Multithreading \\
IPC \\
Optimisation performance et mémoire \\
Réseau \\
Systèmes distribués

\vspace{3mm}

\textbf{OUTILS}

Syslog \\
Docker \\
Git \\
Bash \\
GCC / G++ \\
Make / CMake \\
DBeaver \\
Toolchains cross-compilation \\
FreeCAD \\
PrusaSlicer

\vspace{3mm}
\underline{\textbf{MATÉRIEL}}\\
Prusa XL

\vspace{5mm}

\underline{\textbf{LANGUES}}\\

Français \\
Anglais


% ================= RIGHT COLUMN =================
\switchcolumn

\underline{\textbf{EXPÉRIENCE}}\\

\textbf{Fondateur – Conception hardware et système informatique} \\
2024–2026 (2 ans)

Conception et développement d’un système informatique complet orienté hardware (ordinateur portable sur mesure).  
Travail sur l’architecture matérielle, le choix des composants et la conception système.  
Prototypage et validation via CAO (FreeCAD) et impression 3D (Prusa XL).  
Veille technologique et définition des choix techniques.

\vspace{3mm}

\textbf{Ingénieur développement logiciel C++ – Plateforme de gestion de données distribuées} \\
2022–2023 (18 mois)

- Développement C++ d’une plateforme de gestion et synchronisation de données distribuées  
- Réplication temps réel multi-serveurs  
- Traitement et validation de données (SQLite)  
- Migration Windows → Linux  
- Intégration et déploiement sur IBM i (AS/400)

\vspace{3mm}

\textbf{Ingénieur logiciel embarqué C/C++ – Secteur Défense} \\
2018–2022

- Développement logiciel embarqué critique en C/C++  
- Intégration sur systèmes Linux embarqués  
- Debug bas niveau et analyse système  
- Optimisation performance et robustesse  
- Collaboration équipes hardware et système  
- Gestion d’équipe (3 personnes)

\vspace{5mm}

\underline{\textbf{FORMATIONS}}\\

Licence 2 EEA (Niveau DEUG)

\end{paracol}

\end{document}
